% GUI agent vs API agent
\begin{figure}[H]
\centering
\begin{tikzpicture}
\begin{axis}[
  ybar,
  bar width=18pt,
  ymin=0, ymax=10,
  ylabel={מדד יחסי (גבוה = פחות יעיל)},
  symbolic x coords={GUI Agent, API/Text Agent},
  xtick=data,
  xticklabel style={font=\small},
  legend style={at={(0.5,-0.22)},anchor=north, legend columns=3, font=\footnotesize},
  width=11cm, height=6cm,
  nodes near coords,
  nodes near coords style={font=\tiny}
]
\addplot[fill=accentamber!70] coordinates {(GUI Agent,8) (API/Text Agent,2)};
\addlegendentry{עלות טוקנים}
\addplot[fill=accentblue!60] coordinates {(GUI Agent,9) (API/Text Agent,3)};
\addlegendentry{מורכבות}
\addplot[fill=accentgreen!60] coordinates {(GUI Agent,7) (API/Text Agent,2)};
\addlegendentry{זמן ביצוע}
\end{axis}
\end{tikzpicture}
\caption{סוכן דפדפן מול סוכן טקסט/\eng{API}: עלות, מורכבות וזמן (אילוסטרציה לפי ההרצאה)}
\label{fig:gui-vs-api}
\end{figure}
